\chapter{Đánh giá hệ thống}

\section{Phương pháp và tiêu chí đánh giá}

Đánh giá hệ thống nhằm trả lời năm câu hỏi:
\begin{enumerate}
    \item Hệ thống có đáp ứng các yêu cầu chức năng chính không?
    \item Luồng tích hợp với LMS có hoạt động đúng role/context không?
    \item Retrieval có trả về đúng chunk/card liên quan cho query học tập
          không?
    \item Nội dung do AI sinh ra có đúng kiến thức, rõ ràng và phù hợp
          Bloom/LO không?
    \item Hiệu năng pipeline và API có đủ cho demo/staging không?
\end{enumerate}

Luận văn không tạo số liệu giả. Các số liệu chưa đo đầy đủ được đánh dấu
\texttt{TODO: đo thực nghiệm} hoặc để \ldots{} trong bảng kết quả. Khi báo
cáo kết quả cuối cùng, cần ghi rõ tập dữ liệu, prompt version, model version,
cấu hình retrieval và tiêu chí chấm để đảm bảo khả năng tái lập.

\section{Đánh giá tính năng và tích hợp}

\subsection{Functional Tests}

Kiểm thử chức năng bao phủ các flow chính: upload tài liệu, xử lý pipeline,
review, publish, học bằng card, làm quiz, semantic search, ingest curriculum
và LTI launch.

\begin{longtable}{|p{1.8cm}|p{5cm}|p{4cm}|p{3cm}|}
    \caption{Kịch bản kiểm thử chức năng.} \label{tab:functional-tests-lvtn} \\
    \hline
    \textbf{ID} & \textbf{Kịch bản} & \textbf{Kỳ vọng} & \textbf{Kết quả} \\
    \hline
    \endfirsthead
    \hline
    \textbf{ID} & \textbf{Kịch bản} & \textbf{Kỳ vọng} & \textbf{Kết quả} \\
    \hline
    \endhead

    FT-01 & Upload PDF hợp lệ & Document tạo, job queued & TODO: đo thực nghiệm \\ \hline
    FT-02 & Upload file sai định dạng & Hệ thống từ chối, báo lỗi rõ & TODO: đo thực nghiệm \\ \hline
    FT-03 & Pipeline parse/chunk/embed & Chunk và vector được lưu & TODO: đo thực nghiệm \\ \hline
    FT-04 & Sinh card/quiz & Card/quiz đúng schema & TODO: đo thực nghiệm \\ \hline
    FT-05 & Review và publish & Học viên chỉ thấy nội dung đã publish & TODO: đo thực nghiệm \\ \hline
    FT-06 & Học viên làm quiz & Lưu submission và trả feedback & TODO: đo thực nghiệm \\ \hline
    FT-07 & Semantic search & Trả về top-K nội dung liên quan & TODO: đo thực nghiệm \\ \hline
    FT-08 & LTI launch & Session tạo đúng user/context & TODO: đo thực nghiệm \\ \hline
\end{longtable}

\subsection{Integration Tests}

Các integration test cần chạy với Docker Compose hoặc test container cho
PostgreSQL, Qdrant, Neo4j, Redis và MinIO. Mục tiêu là kiểm tra contract giữa
backend gRPC, worker, queue và các adapter hạ tầng.

\begin{table}[H]
    \centering
    \caption{Nhóm kiểm thử kỹ thuật cần duy trì.}
    \label{tab:test-groups-lvtn}
    \begin{tabular}{p{4cm}p{9.5cm}}
        \toprule
        \textbf{Nhóm test} & \textbf{Mục tiêu} \\
        \midrule
        Domain tests & Kiểm tra entity, validation và rule nghiệp vụ. \\
        Application tests & Kiểm tra use case với mock port. \\
        Adapter tests & Kiểm tra PostgreSQL, Qdrant, Neo4j, parser và LLM adapter. \\
        Delivery tests & Kiểm tra gRPC/HTTP contract và error mapping. \\
        E2E demo tests & Kiểm tra upload đến học/làm quiz trên stack thật. \\
        LTI tests & Kiểm tra OIDC login, launch JWT, role/context mapping và session. \\
        \bottomrule
    \end{tabular}
\end{table}

\section{Đánh giá chất lượng truy xuất}

\subsection{Vector Retrieval}

Vector retrieval được đánh giá bằng một bộ query và ground truth. Mỗi query
được gắn với danh sách chunk hoặc card đúng. Bộ query nên bao gồm query trực
tiếp theo thuật ngữ trong tài liệu, query diễn giải lại bằng ngôn ngữ tự
nhiên, query theo Learning Outcome và query gây nhiễu nằm ngoài phạm vi khóa
học.

\subsection{Graph-aware Retrieval}

Graph-aware retrieval được so sánh với vector-only và vector + metadata
filter. Mục tiêu là kiểm tra liệu quan hệ Course, Chapter, Learning Outcome và
Chunk có giúp tăng khả năng chọn đúng context khi sinh quiz hoặc trả lời câu
hỏi theo LO hay không.

\subsection{Retrieval Accuracy}

\begin{table}[H]
    \centering
    \caption{Chỉ số đánh giá truy xuất.}
    \label{tab:retrieval-metrics-lvtn}
    \begin{tabular}{p{3cm}p{10.5cm}}
        \toprule
        \textbf{Chỉ số} & \textbf{Ý nghĩa} \\
        \midrule
        Recall@K & Tỷ lệ query có ít nhất một chunk đúng trong top-K. \\
        MRR & Đo vị trí trung bình của kết quả đúng đầu tiên. \\
        nDCG@K & Đánh giá thứ hạng khi ground truth có mức liên quan khác nhau. \\
        Latency P95 & Thời gian truy xuất ở phân vị 95. \\
        \bottomrule
    \end{tabular}
\end{table}

\begin{longtable}{|p{3cm}|p{3cm}|p{3cm}|p{4cm}|}
    \caption{Bảng kết quả retrieval dự kiến.} \label{tab:retrieval-results-lvtn} \\
    \hline
    \textbf{Chiến lược} & \textbf{Recall@5} & \textbf{MRR} & \textbf{Ghi chú} \\
    \hline
    \endfirsthead
    \hline
    \textbf{Chiến lược} & \textbf{Recall@5} & \textbf{MRR} & \textbf{Ghi chú} \\
    \hline
    \endhead

    Vector-only & \ldots & \ldots & TODO: đo thực nghiệm \\ \hline
    Vector + metadata filter & \ldots & \ldots & TODO: đo thực nghiệm \\ \hline
    Vector + graph-aware rerank & \ldots & \ldots & TODO: đo thực nghiệm \\ \hline
\end{longtable}

\section{Đánh giá chất lượng nội dung sinh ra}

\subsection{Quiz Quality}

Quiz quality được đánh giá theo các tiêu chí:
\begin{itemize}
    \item \textbf{Đúng kiến thức:} câu hỏi và đáp án đúng với tài liệu gốc.
    \item \textbf{Rõ ràng:} câu hỏi không mơ hồ, không đánh đố vô lý.
    \item \textbf{Phù hợp độ khó:} phù hợp trình độ môn học và Bloom level.
    \item \textbf{Chất lượng distractor:} đáp án sai hợp lý, không quá lộ.
    \item \textbf{Gắn nguồn:} có thể truy ngược về chunk hoặc LO liên quan.
\end{itemize}

Giảng viên hoặc trợ giảng chấm mỗi câu quiz theo thang Likert 1--5. Một câu
được xem là đạt nếu điểm trung bình $\geq$ 3.5 và không có lỗi kiến thức
nghiêm trọng.

\subsection{Flashcard Quality}

Flashcard hoặc lesson card cần được đánh giá ở ba khía cạnh: mức độ đúng với
nguồn, khả năng tóm tắt đủ ý chính và độ phù hợp với học tập ngắn. Card quá
dài làm mất ý nghĩa Micro-learning, trong khi card quá ngắn có thể thiếu ngữ
cảnh. Do đó cần đo cả độ dài, khả năng truy ngược nguồn và phản hồi của người
dùng.

\subsection{Learning Content Quality}

\begin{longtable}{|p{4.5cm}|p{3cm}|p{3cm}|p{3cm}|}
    \caption{Bảng đánh giá chất lượng nội dung sinh ra.} \label{tab:quiz-quality-lvtn} \\
    \hline
    \textbf{Tiêu chí} & \textbf{Điểm TB} & \textbf{Tỷ lệ đạt} & \textbf{Ghi chú} \\
    \hline
    \endfirsthead
    \hline
    \textbf{Tiêu chí} & \textbf{Điểm TB} & \textbf{Tỷ lệ đạt} & \textbf{Ghi chú} \\
    \hline
    \endhead

    Đúng kiến thức & \ldots/5 & \ldots\% & TODO: đo thực nghiệm \\ \hline
    Rõ ràng & \ldots/5 & \ldots\% & TODO: đo thực nghiệm \\ \hline
    Phù hợp Bloom/LO & \ldots/5 & \ldots\% & TODO: đo thực nghiệm \\ \hline
    Chất lượng distractor & \ldots/5 & \ldots\% & TODO: đo thực nghiệm \\ \hline
    Giải thích đáp án & \ldots/5 & \ldots\% & TODO: đo thực nghiệm \\ \hline
    Chất lượng flashcard & \ldots/5 & \ldots\% & TODO: đo thực nghiệm \\ \hline
\end{longtable}

\section{Đánh giá hiệu năng}

\subsection{Response Time}

Response time tập trung vào các API người dùng tương tác trực tiếp như đọc
card, lấy quiz, nộp đáp án và semantic search. Các API đọc card/quiz nên có
mục tiêu P95 $\leq$ 500ms trong môi trường staging; semantic search có thể
chấp nhận cao hơn do cần embedding query và truy xuất vector.

\subsection{Throughput}

Throughput đo khả năng phục vụ đồng thời của backend và worker. Với API, chỉ
số quan trọng là số request/giây ở mức latency chấp nhận được. Với worker, chỉ
số quan trọng là số chunk hoặc document xử lý được trong một khoảng thời gian
với cấu hình tài nguyên cố định.

\subsection{Async Processing Performance}

Hiệu năng xử lý bất đồng bộ cần đo thời gian parse, chunking, embedding,
generation và persist. Các bước gọi LLM thường chiếm phần lớn thời gian nên
cần ghi riêng latency, số lần retry và tỷ lệ lỗi provider.

\begin{longtable}{|p{4cm}|p{3cm}|p{3cm}|p{4cm}|}
    \caption{Bảng kết quả hiệu năng dự kiến.} \label{tab:performance-lvtn} \\
    \hline
    \textbf{Chỉ số} & \textbf{Mục tiêu} & \textbf{Kết quả} & \textbf{Ghi chú} \\
    \hline
    \endfirsthead
    \hline
    \textbf{Chỉ số} & \textbf{Mục tiêu} & \textbf{Kết quả} & \textbf{Ghi chú} \\
    \hline
    \endhead

    Xử lý tài liệu 20 trang & $\leq$ 3 phút & \ldots & TODO: đo thực nghiệm \\ \hline
    API get cards P95 & $\leq$ 500ms & \ldots & TODO: đo thực nghiệm \\ \hline
    API get quiz P95 & $\leq$ 500ms & \ldots & TODO: đo thực nghiệm \\ \hline
    Search P95 & $\leq$ 2s & \ldots & TODO: đo thực nghiệm \\ \hline
    Concurrent users & 50 & \ldots & TODO: đo thực nghiệm \\ \hline
    Worker throughput & \ldots & \ldots & TODO: đo thực nghiệm \\ \hline
\end{longtable}
